%! Transformations of Functions — Unit 1, Lesson 1.4 — Slides (Beamer)
\documentclass[aspectratio=169, 11pt]{beamer}
\usepackage{precalculus-beamer}   % provides \plumheader{} and \sectionlabel[color]{}

\begin{document}

% TITLE SLIDE
{
\setbeamercolor{background canvas}{bg=plum}
\begin{frame}[plain]
  \vspace{2.8cm}\hspace{0.55cm}%
  \begin{minipage}{0.9\textwidth}
    {\color{goldacc}\bfseries\small LESSON 1.4}\\[0.5cm]
    {\color{white}\bfseries\LARGE Transformations of Functions}\\[1.2cm]
    {\color{lilacmid}\small Precalculus~$\cdot$~Unit 1: Functions and Change}
  \end{minipage}
\end{frame}
}

% HOOK SLIDE
\begin{frame}[t, plain]
  \plumheader{Hook: Two Harbors}
  \sectionlabel[goldacc]{Vote Before You Compute}

  \(D(t)\) is the water depth at North Harbor, \(t\) hours after midnight.
  South Harbor rises and falls the same way, and its depth is published as
  \(S(t) = D(t + 3)\).

  \smallskip
  \begin{columns}[T, onlytextwidth]
    \column{0.54\textwidth}
      \begin{center}
      \renewcommand{\arraystretch}{1.3}
      \begin{tabular}{|l|c|c|c|c|c|}
      \hline
      \(t\) (hours) & 0 & 3 & 6 & 9 & 12 \\
      \hline
      \(D(t)\) (feet) & 4 & 8 & 12 & 8 & 4 \\
      \hline
      \end{tabular}
      \end{center}

      \medskip
      The rule \textbf{adds} three hours.

    \column{0.42\textwidth}
      \begin{block}{The question}
        Does South Harbor's high tide come \textbf{before} or \textbf{after}
        North Harbor's?
      \end{block}

      \medskip
      \begin{block}{Then compute}
        \(S(3) = D(6) = 12\). South peaks at 3 a.m., three hours
        \textbf{ahead}.
      \end{block}
  \end{columns}
\end{frame}

% SECTION 1 — QUESTION AND ANSWER
\begin{frame}[t, plain]
  \plumheader{1. The Question and the Answer}
  \sectionlabel[redacc]{The Whole Lesson in One Idea}

  Using a function is a transaction: you hand \(f\) a question, \(f\) hands
  back an answer. A rule can interfere in exactly two places.

  \medskip
  \begin{columns}[T, onlytextwidth]
    \column{0.44\textwidth}
      \begin{center}
      \begin{tikzpicture}
        \node[font=\large, anchor=west] (gq) at (0,0) {\(g(x) =\)};
        \node[font=\large, goldacc, anchor=west] (a) at (gq.east) {\(\;a\)};
        \node[font=\large, anchor=west] (mid) at (a.east) {\(\cdot\, f(\)};
        \node[font=\large, plum, anchor=west] (inner) at (mid.east) {\(\,x + h\,\)};
        \node[font=\large, anchor=west] (close) at (inner.east) {\()\)};
        \node[font=\large, goldacc, anchor=west] (k) at (close.east) {\(\;+\;k\)};
        \node[font=\scriptsize\bfseries, plum, align=center]
          (ql) at ([yshift=-0.9cm]inner.south) {inside: the question};
        \draw[thick, plum, ->] (ql.north) -- (inner.south);
        \node[font=\scriptsize\bfseries, goldacc, align=center]
          (al) at ([yshift=0.8cm]mid.north) {outside: the answer};
        \draw[thick, goldacc, ->] (al.west) to[out=180, in=105] (a.north);
        \draw[thick, goldacc, ->] (al.east) to[out=0, in=80] (k.north);
      \end{tikzpicture}
      \end{center}

    \column{0.52\textwidth}
      \begin{block}{Outside \(\Rightarrow\) the answer}
        \textbf{Vertical}, and it does exactly what it says.
      \end{block}

      \medskip
      \begin{block}{Inside \(\Rightarrow\) the question}
        \textbf{Horizontal}, and it runs \textbf{opposite} to what it looks
        like --- because you must work out which \(x\) makes the question come
        out right.
      \end{block}
  \end{columns}
\end{frame}

% SECTION 2 — VERTICAL
\begin{frame}[t, plain]
  \plumheader{2. Rewriting the Answer --- Vertical Moves}

  \begin{columns}[T, onlytextwidth]
    \column{0.44\textwidth}
      \begin{center}
      \begin{tikzpicture}[x=0.66cm, y=0.24cm]
        \draw[linegray, very thin] (-3,-5) grid (3,8);
        \draw[->, thick] (-3.4,0) -- (3.6,0);
        \draw[->, thick] (0,-5.5) -- (0,8.7);
        \foreach \y in {-4,4,8}
          \node[left, font=\tiny, fill=white, inner sep=1pt] at (0,\y) {\y};
        \draw[thick, goldacc, dashed, smooth, samples=60, domain=-2.2:2.2]
          plot (\x, {\x*\x + 3});
        \draw[thick, slate, dashed, smooth, samples=60, domain=-2.2:2.2]
          plot (\x, {-\x*\x});
        \draw[very thick, plum, smooth, samples=60, domain=-2.7:2.7]
          plot (\x, {\x*\x});
        \node[font=\tiny, goldacc, fill=white, inner sep=1pt] at (-1.6,7.4) {\(f(x)+3\)};
        \node[font=\tiny, plum, fill=white, inner sep=1pt] at (1.8,7.4) {\(f(x)\)};
        \node[font=\tiny, slate, fill=white, inner sep=1pt] at (1.7,-4.3) {\(-f(x)\)};
      \end{tikzpicture}
      \end{center}

    \column{0.52\textwidth}
      \begin{itemize}\setlength{\itemsep}{6pt}
        \item \(f(x) + k\) --- \textbf{translation}: every answer gains \(k\).
              Up if \(k > 0\), down if \(k < 0\).
        \item \(af(x)\) --- \textbf{dilation} by \(|a|\): stretched away from
              the \(x\)-axis if \(|a| > 1\), compressed toward it if
              \(0 < |a| < 1\).
        \item \(a < 0\) also carries a \textbf{reflection} over the
              \(x\)-axis.
      \end{itemize}

      \medskip
      \begin{block}{The consequence}
        The questions were never touched --- so the \textbf{domain} is
        unchanged and the \textbf{range} moves.
      \end{block}
  \end{columns}
\end{frame}

% SECTION 3 — HORIZONTAL
\begin{frame}[t, plain]
  \plumheader{3. Rewriting the Question --- Horizontal Moves}
  \sectionlabel[goldacc]{The One Students Get Backwards}

  \begin{columns}[T, onlytextwidth]
    \column{0.50\textwidth}
      \(g(x) = f(x + 3)\), with \(f(x) = x^2\):

      \smallskip
      \begin{center}
      \renewcommand{\arraystretch}{1.3}
      \begin{tabular}{|l|c|c|c|c|c|}
      \hline
      \(x\) & \(-5\) & \(-4\) & \(-3\) & \(-2\) & \(-1\) \\
      \hline
      handed to \(f\) & \(-2\) & \(-1\) & 0 & 1 & 2 \\
      \hline
      \(g(x)\) & 4 & 1 & 0 & 1 & 4 \\
      \hline
      \end{tabular}
      \end{center}

      \medskip
      The output list is \textbf{identical}. Only the addresses moved.

    \column{0.46\textwidth}
      \begin{block}{Where did the vertex go?}
        \(f\) was handed \(0\) at its vertex, so solve
        \[ x + 3 = 0 \;\Longrightarrow\; x = -3. \]
        \textbf{Left} 3 --- not the direction the ``\(+\)'' suggests.
      \end{block}

      \medskip
      \(f(2x)\): to hand \(f\) its \(2\), supply \(x = 1\). Every landmark
      arrives at \textbf{half} the \(x\) --- a squeeze by \(1/2\).

      \smallskip
      Questions changed, answers untouched: the \textbf{range} is unchanged
      and the \textbf{domain} moves.
  \end{columns}
\end{frame}

% SECTION 4 — FOLLOW THE LANDMARK
\begin{frame}[t, plain]
  \plumheader{4. Follow the Landmark}
  \sectionlabel[redacc]{No Order of Operations to Memorize}

  \begin{block}{For \(g(x) = a\,f(x+h) + k\), the parent's point \((p, q)\) lands at}
    \[
      \bigl(\,p - h,\;\; aq + k\,\bigr)
      \qquad\text{new input: solve } x + h = p.
    \]
  \end{block}

  \medskip
  \(g(x) = -2f(x+1) - 3\), with \(f(x) = x^2\):

  \smallskip
  \begin{center}
  \renewcommand{\arraystretch}{1.35}
  \begin{tabular}{|l|c|c|c|}
  \hline
  \textbf{parent point} & \textbf{new input} \(p - 1\) & \textbf{new output} \(-2q - 3\) & \textbf{lands at} \\
  \hline
  \((-2,\,4)\) & \(-3\) & \(-11\) & \((-3,\,-11)\) \\
  \hline
  \((0,\,0)\) & \(-1\) & \(-3\) & \((-1,\,-3)\) \\
  \hline
  \((2,\,4)\) & 1 & \(-11\) & \((1,\,-11)\) \\
  \hline
  \end{tabular}
  \end{center}

  \medskip
  Left 1, stretched by 2, flipped over the \(x\)-axis, down 3 --- and the
  vertex, once the lowest point, is now the \textbf{highest}.
\end{frame}

% SECTION 5 — DOMAIN AND RANGE
\begin{frame}[t, plain]
  \plumheader{5. The Two Columns Carry Domain and Range}

  The domain is a fact about the \textbf{input} column; the range is a fact
  about the \textbf{output} column.

  \medskip
  \begin{columns}[T, onlytextwidth]
    \column{0.48\textwidth}
      \begin{block}{Domain --- from the question}
        \(f\) has domain \([-2, 4]\), and \(g(x) = -2f(x+1) - 3\):
        \[ -2 \le x + 1 \le 4 \;\Longrightarrow\; -3 \le x \le 3 \]
      \end{block}

    \column{0.48\textwidth}
      \begin{block}{Range --- from the answer}
        \(f\) has range \([0, 9]\), and every \(q\) becomes \(-2q - 3\):
        \[ -2(0) - 3 = -3, \qquad -2(9) - 3 = -21 \]
        Range \([-21, -3]\) --- the endpoints \textbf{trade places}.
      \end{block}
  \end{columns}

  \medskip
  \begin{center}
    \textbf{Changes to the question move the domain. Changes to the answer
    move the range.}
  \end{center}
\end{frame}

% CLASS PRACTICE
\begin{frame}[t, plain]
  \plumheader{Class Practice}
  \sectionlabel[redacc]{Try It}

  \begin{enumerate}\setlength{\itemsep}{10pt}
    \item A function \(f\) has this table, with its peak at \((2, 4)\):

          \smallskip
          \begin{tabular}{|l|c|c|c|c|c|}
          \hline
          \(x\) & 0 & 1 & 2 & 3 & 4 \\
          \hline
          \(f(x)\) & 1 & 3 & 4 & 3 & 1 \\
          \hline
          \end{tabular}

          \smallskip
          Fill the row for \(g(x) = f(x) + 2\). Then let \(h(x) = f(x-2)\):
          what is \(h(5)\), and where is \(h\)'s peak?

    \item The lowest point on the graph of \(f\) is \((3, -1)\) and its domain
          is \([1, 7]\). For \(g(x) = 4f(x-2) + 5\), find the lowest point and
          the domain --- and say which move changed the domain.
  \end{enumerate}
\end{frame}

\end{document}
